
\section{Causal rising covers and the full native upper bound}
\label{upper:section}

This section proves the upper bound on the original divisor union.
The main step treats a proper full-to-zero native transition with two
small endpoint reserves.  A conditional rising cover supplies both exponential half-line kernels
needed by the original routing filtration.  Their powers are the
products of their respective owners, rather than the final product.

\subsection{Notation and previously established inputs}
\label{upper:inputs-section}

Fix the dimension bound $H$.  Perform the fixed bounded lift once and
take its genuine ordered optimizer.  Its maximal equal-product owners
have products
\[
  0<\sigma_*\leq \sigma_m<\cdots<\sigma_1\leq\sigma^*<1.
\]
All constants below depend on $H$ and the fixed conventions in the
previously established estimates.  Write
\begin{equation}
 s=\sigma_1,\qquad r_j=\frac{\sigma_j}{s},\qquad
 u_j=\frac{\sigma_{j+1}}{\sigma_j},\qquad
 d_j=\frac1{u_j}-1\quad(j<m).
 \label{upper:ratios}
\end{equation}
The compound child weight of a proper owner is denoted by $z_j$.
It satisfies $1<z_*\leq z_j\leq H$, and a native slot with $r$
eligible exterior names has original coarse law
\begin{equation}
 \pi_j(c)=\frac1{r+z_j}\quad(c\text{ exterior}),\qquad
 \pi_j(\mathrm{child})=\frac{z_j}{r+z_j}.
 \label{upper:native-law}
\end{equation}

The following are the precise earlier inputs used in the proof.
Their complete proofs are given in Appendix~\ref{app:upper-details};
we state them together to make the parameters of the conditional
cover argument explicit.

\begin{proposition}[Geometric and scalar estimates]
\label{upper:established-inputs}
The following assertions hold with constants depending only on $H$.
\begin{enumerate}
\item Every retained native transition is canonical
$r\to q$, with full input and deepest-suffix output.  Within an owner,
sufficiently rare transitions have disjoint rank intervals: if one
precedes another, then $q_1\geq r_2$.  Thus a zero-output transition,
if present, is last.  Its actual intrinsic endpoint scores are positive,
and its persistence factor is
\[
 \Psi_i=\min\left\{1,
 \left(\frac{(1+X_i)(1+Y_i)}{1+\mu_i}\right)^{2\chi_i}\right\}.
\]
The exponents and rates are in fixed compact ranges.  Factors above
any fixed positive rarity threshold are bounded below by a fixed
positive constant.

\item A retained nonzero-output transition has a necessary test on
an original short native bank.  The test reads that bank's coarse
labels and the uncolored counts in its intermediate native groups.
The corresponding true failure cover has volume at most $S_i e^{-h}$.
Its private probability $p_i(h)$ satisfies
\[
 \mathbb E p_i(h)^{\sigma_{o(i)}}
       \leq C(1+h)^K\Psi_i .
\]
The same assertion holds for a proper zero-output transition when
$1+X_i\geq c_X\sqrt{1+L_i}$.  In that case its intermediate groups
are all the later native groups of its owner.  These current-plus-
intermediate groups are pairwise disjoint within each owner and across
owners.  All these estimates allow fixed point insertions at the same
power.  Original count middles are disjoint from the short label banks.

\item For any proper full-to-zero transition which remains genuinely
rare, let $M$ be its native block length, $A,B$ the preceding and
following native lengths of that owner, and $V_{\rm old}$ its complete
old physical length.  Then
\begin{align}
 V_{\rm old}+A+B&\leq C\tau M,\qquad T_j\asymp M,
 \label{upper:dominance}\\
 L&=1+B,\qquad R=1+V_{\rm old}+A+\varrho,\notag\\
 \varrho&=\max\left\{\varrho_0,
       \frac{c_0\sqrt{1+M}}{1+d_j\sqrt{1+M}}\right\},\notag\\
 X_i&\asymp R,\qquad Y_i\asymp L,\qquad
 |e|M\leq C(V_{\rm old}+A+B),
 \label{upper:zero-scales}
\end{align}
where $e=(r+z_j)^{\sigma_j}\mathbb EW-r$ and
$W=\log(r+z_j)-\log z_j\,1_{\rm child}$.
The $V_{\rm old}$ here consists of the genuine original groups and
their transported coefficients.  It is not a newly sampled old list.

\item On disjoint original end banks and a middle bank in the $M$
block, the signed logarithmically lowered one-end estimates and the
whole-band binomial maximum atom give
\begin{equation}
 \mathbb E B_w(h,n)^{\sigma_j}
 \leq C(1+h)^K(1+|n|)^K F_j^0,
 \qquad
 F_j^0=\left(\frac{\varrho}{\sqrt{1+M}}\right)^{\sigma_j}
       \min\left\{1,\left(\frac{LR}{1+M}\right)^{2\chi_i}\right\}.
 \label{upper:scalar-bank}
\end{equation}
Here $w\geq M/2$; a completed-word band has width $\varrho$, while a
first-hit band may have fixed width.  The left reserve is at most
$h+C(1+N_B)$ and the right reserve is at most
$C(1+h+V_{\rm old}+A+(|n|+1)\varrho)$.  The estimates sum all
original counts and keep all $n+1$ conditional gaps.  If the larger
reserve is superdiffusive, only the smaller-reserve bank of length
$cM^2/\max(L,R)^2$ is retained.  In the small-gap case
$d_j\sqrt M<1$, a left bank alone gives the same factor, since
$\varrho/\sqrt M\asymp1$.

\item Fixed-dimensional original shot sums have all fixed moments.
Their insertion estimates are obtained by a finite Poisson--Mecke
partition before summing bank locations.  A count middle occurring in
an intermediate reserve group is handled by the conditional inequality
\begin{equation}
 \mathbb E[f(N)g(N)]\leq\mathbb Ef(N)\,\mathbb Eg(N)
 \quad(f\text{ nonincreasing},\ g\text{ nondecreasing}),
 \label{upper:opposite-monotonicity}
\end{equation}
with all other groups fixed.  This includes its count weight
$(1+d_j\sqrt{1+N})^{-1}$ and an increasing necessary-bank reserve.
\end{enumerate}
\end{proposition}

For clarity, the bounded-scaled-drift hypothesis in
\eqref{upper:scalar-bank} follows from \eqref{upper:zero-scales}:
use two banks of length $M/16$ if $\max(L,R)\leq C\sqrt M$,
and otherwise a bank of length $cM^2/\max(L,R)^2$.  The original
correlation parameter and the centered scalar parameter differ by
$O(LR/M)$, so the exponent comparison costs only
$\exp\{O((LR/M)|\log(LR/M)|)\}$.  Thus \eqref{upper:scalar-bank}
has the same current intrinsic factor as the lower bound.

Let $F_j$ be the proper owner's count factor, at power $\sigma_j$,
times all its retained native persistence factors.  For the terminal
owner it is just the product of its retained nondegenerate factors.
Consequently
\begin{equation}
 \prod_{j=1}^m F_j=\mathcal C_{\rm count}\mathcal P_{\rm ret}.
 \label{upper:factor-product}
\end{equation}
Each $F_j$ has a dimension-dependent polynomial lower bound in the
total length of the original native groups it reads.  The terminal
ordinary factor is denoted by $\mathcal B_I$ and is kept separate.

\subsection{Actual routing capacities and two left covers}
\label{upper:geometry-section}

By Lemma~\ref{artrunc:lemma}, the full squarefree harmonic volume sum is
at most twice its restriction to $a_i\leq y_i^{\tau_k}$ for all $i$.
The prefix bounds \eqref{artrunc:prefix-size} supply the dimension-constant
full-size rectangles below. Fix such an original squarefree tuple and its complete
uncolored data $G$.  Its original product fine-label measure is
$\pi_G$.  Let $\mathcal Q_j$ reveal the complete first $j$ owner
routes and leave the genuine descendant alphabet compounded.
Write $Z_j$ for the actual grouped output and put
\begin{equation}
 P_j=\sum_{\ell\leq j}Z_\ell,\qquad
 \eta_j=\sum_{\ell\leq j}\frac{\sigma_\ell}{\sigma_j}Z_\ell,
 \qquad X=\sum_{j=1}^m r_j Z_j.
 \label{upper:outputs}
\end{equation}
The common root $X$ is uncolored.  For $j<m$ the exact identities are
\begin{equation}
 P_j=\eta_j-\sum_{\ell<j}d_\ell\eta_\ell,
 \qquad
 P_m=\frac{X}{r_m}-\sum_{\ell<m}d_\ell\eta_\ell.
 \label{upper:carry-identities}
\end{equation}

Call a proper owner \emph{hard} if its retained zero-output edge does
not satisfy the one-end admissibility hypothesis of
Proposition~\ref{upper:established-inputs}(2).  All other owners are
simple.  A hard owner has one selected zero block and all its other
retained edges are earlier nonzero edges.  Constants may be chosen
so that a hard owner has $d_j\sqrt M\geq1$; the small-gap variant below
also covers the bounded transition range.

\begin{lemma}[Original projected rectangles and dual left walls]
\label{upper:dual-walls}
A complete $\mathcal Q_j$ atom $q$ has actual capacity
\begin{equation}
 \operatorname{Vol}(q)\leq C\pi_G(q)e^{P_j(q)}.
 \label{upper:complete-cap}
\end{equation}
At a hard owner, fix a genuine $\mathcal Q_{j-1}$ atom $\alpha$.
Read its selected zero block backwards and set
\begin{equation}
 C_j(w)=H_B+\sum_{x<w}
       [\log(r+z_j)-\log z_j\,1_{\rm child}]-rw,
 \qquad Z_j=C_j(M)+Y_A.
 \label{upper:local-risk}
\end{equation}
Then $H_B\leq C N_B$ and
$Y_A\geq-C(V_{\rm old}+A)$ wordwise.  There are two genuine
partial capacities,
\begin{align}
 \operatorname{Vol}(Q_w^{\rm glob})
   &\leq C\pi_G(Q_w^{\rm glob})e^{C_j(w)}\operatorname{Shot}_w,
 \label{upper:global-native-cap}\\
 \operatorname{Vol}(\alpha,Q_w)
   &\leq C\pi_G(\alpha,Q_w)e^{P_{j-1}(\alpha)+C_j(w)}
        \operatorname{Shot}_w.
 \label{upper:cumulative-native-cap}
\end{align}
Imposing both left walls
\begin{equation}
 C_j(w)\geq-h-\mathcal L(w),\qquad
 P_{j-1}+C_j(w)\geq-h-\mathcal L(w),\qquad
 \mathcal L(w)=A_0\log(1+\min(w,M-w)),
 \label{upper:two-left-tests}
\end{equation}
costs at most $CS e^{-h}$ in the union of their failures.  Their
simultaneous survivors have effective reserve
\begin{equation}
 h_{\rm eff}=h+\min(P_{j-1},0)\leq h.
 \label{upper:effective-reserve}
\end{equation}
Here $S\geq1$ is an uncolored original shot with all fixed moments;
it may be chosen jointly for the finitely many owners and other
native failure covers, without reading the final unread root suffix.
\end{lemma}

\begin{proof}
Fixing the observed exterior divisors gives translated original
rectangles in their coordinates.  All unresolved coordinates keep
their original size bounds.  The negative logarithm of the observed
class probability is the sum of its original information marks.
Subtracting the geometric budgets gives precisely $P_j$, proving
\eqref{upper:complete-cap}.

For \eqref{upper:global-native-cap}, leave all ancestor and earlier
coordinates in their full rectangles.  Reveal the later native word
and the crossed points of the current block.  The $r$ unresolved
exterior widths have common threshold $b-w$, where $b$ is the block's
original right endpoint.  Their budget deficit is
$rw+\sum_{\rm later}r_gL_g$.  The negative observed log mass is the
corresponding mark sum, so their difference is $C_j(w)$.  Every other
original coordinate budget cancels.  In particular no earlier private
output or selected incoming count occurs in this capacity.

Instead first fix $\alpha$.  Its actual resolved translations give
$\pi_G(\alpha)e^{P_{j-1}}$.  The same rectangle on its literal child
fiber gives the conditional factor
$\pi_G(Q_w\mid\alpha)e^{C_j(w)}\operatorname{Shot}_w$.
Their product is \eqref{upper:cumulative-native-cap}.  The fiber is
not renormalized by a survival event.  A shot using all original
points dominates its selected-list shot by positivity.  The two
capacities are different covers, not two factorizations of one
artificially independent route.

At the first failure of either wall, including $w=0$, substitute
the failed inequality into its corresponding capacity.  The result
is the original class probability times
$e^{-h-\mathcal L(w)}\operatorname{Shot}_w$.  The first-failure
classes form an antichain in their own reveal order.  For the second
cover, first sum conditional Kraft masses and then their original
ancestor probabilities.  Their total is at most one.  Summing the
fixed logarithmic shot weights gives $CS e^{-h}$ for each cover.
The two reveal orders need not be one filtration: the union bound
retains all their failures.  Finally their intersection is exactly
\eqref{upper:effective-reserve}.  The wordwise bounds in the statement
follow from positive information marks and the original physical
rank budgets; no mean of an old list is used.
\end{proof}

\subsection{Conditional half-line kernels on true incoming fibers}
\label{upper:kernels-section}

We first record the simple-owner calculation.  A proper native slot is
after all ancestor exterior deadlines, so its ancestor category is
deterministically child.  Thus conditioning on any $\mathcal Q_{j-1}$
atom leaves its native labels with exactly the product law
\eqref{upper:native-law}.

\begin{lemma}[Simple-owner kernels]
\label{upper:simple-kernels}
Choose an original native count middle of $n$ slots, disjoint from
the owner's necessary label banks.  Let $E_j(h)$ be their joint event
and $p_j(h,G)$ its private probability.  For every uncolored threshold
$\xi_j$, set $\zeta_j=\eta_j-\xi_j$.  Uniformly in every genuine
$\mathcal Q_{j-1}$ atom,
\begin{align}
 \mathbb E_{\pi_G}[e^{\zeta_j}1_{\{\zeta_j\leq0\}}1_{E_j(h)}
                   \mid\mathcal Q_{j-1}]&\leq Cb_j(h),\notag\\
 \mathbb E_{\pi_G}[e^{-d_j\zeta_j}1_{\{\zeta_j>0\}}1_{E_j(h)}
                   \mid\mathcal Q_{j-1}]&\leq Cb_j(h),
 \label{upper:simple-half-lines}\\
 b_j(h)&=(1+d_j\sqrt{1+n})^{-1}p_j(h,G).\notag
\end{align}
For a terminal owner the corresponding conditional mask mass is bounded
by $b_m(h)=p_m(h,G)$.
\end{lemma}

\begin{proof}
Fix the actual ancestor atom and then the entire current coarse route
outside the middle.  Its masks are now fixed, and
\[
 \zeta_j=H_0-\log z_j\,K,\qquad
 K\sim\operatorname{Bin}\left(n,\frac{z_j}{r+z_j}\right),
\]
where $H_0$ may depend arbitrarily on the old list and outside word.
The binomial maximum atom is at most $C(1+n)^{-1/2}$.  Summing the
negative exponential half-lattice gives $C(1+n)^{-1/2}$; summing
the positive half-lattice gives $C/[d_j\sqrt{1+n}]$, capped by one.
At $d_j=0$ use only the probability bound.  The resulting common
bound is $C(1+d_j\sqrt{1+n})^{-1}$, uniformly in $H_0$.
Release the outside labels only after taking this uniform bound.
Their mask mass is $p_j(h,G)$ on every ancestor fiber.  The terminal
assertion is the same native-only observation without a count weight.
Future child words were not fixed while varying $K$.
\end{proof}

For a hard owner write $P=V_{\rm old}+A$ and let $E_<(h)$ denote
its earlier own native masks.  Their probability $p_h(G_A)$ reads
only original native groups strictly before the selected block.
Choose a fixed $D_A$ and put
\begin{equation}
 \alpha_h=\min\{1,(2+A)^{-D_A}+p_h(G_A)\}.
 \label{upper:short-prefix-floor}
\end{equation}
By its polynomial factor floor, $D_A$ can be chosen so that
\begin{equation}
 p_h\leq\alpha_h\leq1,\quad
 -\log\alpha_h\leq D_A\log(2+A)\leq C(1+P),\quad
 \mathbb E\alpha_h^{\sigma_j}\leq C(1+h)^K\mathcal P_<.
 \label{upper:alpha-bounds}
\end{equation}
An empty earlier collection has $\alpha_h=1$.

\begin{lemma}[Hard-owner kernels and right stops]
\label{upper:hard-kernels}
On a fixed genuine ancestor atom let
\begin{equation}
 x_j=\xi_j-\frac{\sigma_{j-1}}{\sigma_j}\eta_{j-1}
 \quad(j>1),\qquad x_1=\xi_1,
 \qquad \zeta_j=Z_j-x_j.
 \label{upper:local-threshold}
\end{equation}
In addition to the dual left walls, use the right wall
\begin{equation}
 C_j(w)\geq x_j+\log\alpha_h-\log S
                         -A_1\log(1+M-w).
 \label{upper:right-wall}
\end{equation}
Assume its deterministic gate places right-governed first hits after
$M/2$.  There is one uncolored coefficient $b_j(h)$, independent of
the ancestor private word and of $x_j$, with all three properties:
\begin{enumerate}
\item the total right-stop capacity on the ancestor fiber is at most
$C\pi_G(\alpha)e^{P_{j-1}+x_j}b_j(h)$;
\item complete surviving owner words with $\zeta_j\leq0$ have conditional
weighted mass at most $Cb_j(h)$ under weight $e^{\zeta_j}$;
\item complete surviving owner words with $\zeta_j>0$ have conditional
weighted mass at most $Cb_j(h)$ under weight $e^{-d_j\zeta_j}$.
\end{enumerate}
The complete-survivor conclusions themselves do not require the gate.
The small-gap variant omits the right wall and internal right stops.
\end{lemma}

\begin{proof}
First consider complete survivors.  The exact numerical identity
\begin{equation}
 e^\zeta1_{\{\zeta\leq0\}}+e^{-d_j\zeta}1_{\{\zeta>0\}}
 =e^{-x_j}\min\{e^{Z_j},e^{x_j/u_j-d_jZ_j}\}
 \label{upper:numerical-min}
\end{equation}
is only a test-function identity.  No child volume bound is used.
Set $\bar x=x_j-\log S$ and divide the WHOLE output $Z_j-\bar x$
into bands $[n\varrho,(n+1)\varrho)$, $n\in\mathbb Z$.
The right side of \eqref{upper:numerical-min} is at most $S^{d_j}w_n$,
where
\begin{equation}
 w_n=e^{-d_jn\varrho}\ (n\geq0),\qquad
 w_n=e^{-(|n|-1)\varrho}\ (n\leq-1).
 \label{upper:signed-band-weights}
\end{equation}
For positive bands use the second exponential branch and for negative
bands the first.  This remains valid if a shifted band crosses
$\zeta=0$, since the minimum is at most either branch everywhere.
If $d_j\sqrt M\geq1$, both series of $w_n^{\sigma_j}$ times every
fixed polynomial in $|n|$ converge uniformly.

The reflected necessary reserve follows by subtracting
\eqref{upper:right-wall} from $Z_j=C_j(M)+Y_A$:
\begin{align}
 C+Z_j-\bar x-Y_A-\log\alpha_h
 &\leq C_H(1+V_{\rm old}+A)+(|n|+1)\varrho
                                  +D_A\log(2+A)\notag\\
 &\leq C_D(1+V_{\rm old}+A)+(|n|+1)\varrho.
 \label{upper:reflected-reserve}
\end{align}
Both the supplied root level and $\log S$ cancel.  The small left
reserve is still $h+C(1+N_B)$: only for this necessary test we enlarge
$h_{\rm eff}$ to $h$.

For a fixed band, fix the complete current old and earlier native word,
the later native word, and all selected-block labels outside its middle
bank.  The whole-band center is now fixed.  The middle child-count
maximum atom is uniform in that center.  The two end tests use disjoint
original slots and \eqref{upper:reflected-reserve} no longer reads the
earlier private word.  Release the earlier own word with its event
$E_<(h)$, giving its exact mass $p_h\leq\alpha_h$.  Release the true
old word on the ancestor fiber with total mass one.  The later word is
released with its original mass, using $H_B\leq CN_B$ only in the
left necessary reserve.  Thus the conditional mass of the band is at
most $C\alpha_h B_M(h,n)$, uniformly over the actual incoming list.
With
\[
 T_{\rm full}(h)=\sum_{n\in\mathbb Z}w_n B_M(h,n),
\]
each of the two complete half-lines is at most
$CS^{d_j}\alpha_h T_{\rm full}(h)$.

At a right first hit, positivity of the jumps and bounded downward
speed put the risk in a bounded strip below its wall.  This does not
require a bound on a tied positive prime batch.  Substitution into
\eqref{upper:cumulative-native-cap}, divided by
$\pi_G(\alpha)e^{P_{j-1}+x_j}$, gives
\begin{equation}
 C\frac{\alpha_h}{S}(1+M-w)^{-A_1}
        \operatorname{Shot}_w
 \quad\text{times the true conditional prefix mass}.
 \label{upper:right-node-weight}
\end{equation}
The factor $\alpha_h$ comes from the actual capacity, not from
multiplying an unobserved earlier probability into a node.  Prior
survival and the first-hit strip imply the fixed-$w$ end-bank tests
and one fixed-width middle atom.  Their centers can read the fixed
ancestor data, but their upper bounds are uniform.  Let
$T_{\rm right}(h)$ sum these bounds with the weights and shots in
\eqref{upper:right-node-weight}.  One may therefore take
\begin{equation}
 b_j(h)=C\alpha_h\bigl[S^{d_j}T_{\rm full}(h)
                         +S^{-1}T_{\rm right}(h)\bigr].
 \label{upper:hard-cost}
\end{equation}
This proves the three separate statements.

If $d_j\sqrt M<1$, then $\varrho\asymp\sqrt M$ and $R\geq c\sqrt M$.
Keep only the necessary left bank and omit the right wall.  The function
in \eqref{upper:numerical-min} is at most one.  Whole-word conditioning
gives $p_h$ times this bank's probability; its moment is $F_j^0$ because
the count-width factor is comparable to one.  This avoids any series
with a decay rate tending to zero.
\end{proof}

\begin{lemma}[Joint moments at all owner powers]
\label{upper:joint-moments}
The simple and hard costs can be chosen so that, for every subset
$J\subseteq\{1,\ldots,m\}$, every nonnegative integer vector $(h_j)$,
and every fixed $p$,
\begin{align}
 \mathbb E\left[S^p\prod_{j\in J}b_j(h_j)^{\sigma_j}\right]
 &\leq C_p\prod_{j\in J}(1+h_j)^K F_j,
 \label{upper:subset-moment}\\
 \mathbb E\left[(1+N_{\rm pre})S^p
                  \prod_{j\in J}b_j(h_j)^{\sigma_j}\right]
 &\leq C_p(1+V_{\rm pre})
                      \prod_{j\in J}(1+h_j)^K F_j.
 \label{upper:subset-count-moment}
\end{align}
Moreover $b_j(h)\leq CS^Q$ pointwise, uniformly in $h$, for a fixed $Q$.
These are expectations over the ORIGINAL uncolored Poisson restrictions,
not over clocks conditioned on a chosen stopping route.
\end{lemma}

\begin{proof}
For a hard owner, before attaching shots the input of $\alpha_h$ is
strictly before its selected block.  The fixed-$w$, fixed-band bank
expression reads only that block and later native counts.  Original
Poisson restrictions on these groups are independent.  Apply
\eqref{upper:alpha-bounds} and \eqref{upper:scalar-bank} at the same
power $\sigma_j$, then sum the signed bands and weighted positions.
This gives $\mathcal P_<F_j^0=F_j$ with a fixed polynomial in $h$.

For a simple owner, a middle count in an intermediate group is not
independent of the corresponding mask reserve.  Condition on other
groups and use \eqref{upper:opposite-monotonicity} once, then integrate
the disjoint groups.  The same argument includes an admissible
large-left zero edge.  No input group leaves its own native owner.
Different owners' resulting functions therefore use disjoint original
Poisson restrictions.  Old physical lengths in a hard reserve are
deterministic and introduce no old random count into its cost.

Attach a global shot by a finite Mecke partition to each fixed bank
expression before summing positions or forming any reserve envelope.
Inserted native points are deleted at the cost of a bounded reserve
shift; a later count raises $N_B$ by a bounded amount; an inserted
middle point decreases its inverse-count bound.  The earlier mask has
$\alpha_h(G+\text{points})\leq C_k\alpha_{h+Ck}(G)$, with unchanged
floor.  Apply the preceding disjoint-group estimates after these
fixed insertions.  This proves \eqref{upper:subset-moment}; a raw-count
insertion costs its own mean and proves \eqref{upper:subset-count-moment}.
Omitted owners have factor one, so the proof holds for every subset.
In particular it applies to mixed terms created by deterministic floors.
There is no redistribution of powers by Holder's inequality.

Finally every private bank mass and truncated maximum-atom upper is at
most one, the signed band weights have uniformly bounded sum, and the
weighted right-position sum is bounded by $S$.  The exponents $d_j$
are bounded.  This proves the pointwise bound, which will pay positive
logarithms; it is not inferred from a fractional moment.
\end{proof}

\subsection{Floors, the causal gate, and the original-route tower}
\label{upper:composition-section}

For a proper owner let $T_j$ be its total native length.  For the
terminal owner use only its native head before the unread suffix.
Choose fixed $D_j$ so that $f_j=(2+T_j)^{-D_j}$ satisfies
$f_j^{\sigma_j}\leq CF_j$.  Set
\begin{equation}
 a_j=f_j+\sum_{h=0}^\infty(1+h)^{-\beta}b_j(h),\qquad
 B_* =\sum_{j=1}^m r_j\log a_j,
 \theta_j=\frac1{r_j}\sum_{\ell>j}r_\ell\log a_\ell,
 \qquad \xi_j=\frac X{r_j}+\theta_j.
 \label{upper:pooling}
\end{equation}
The $a_j$ are positive uncolored numbers and need not be at most one.
They are fixed before the single cover reserve $h$ is chosen.  In
particular
\begin{equation}
 b_j(h)\leq a_j(1+h)^\beta,
 \qquad-\log a_j\leq D_j\log(2+T_j).
 \label{upper:pool-bounds}
\end{equation}
The pointwise bound in Lemma~\ref{upper:joint-moments} makes every
sum finite when $\beta>1$.

\begin{lemma}[Mixed pooling and the no-future-scale gate]
\label{upper:pool-gate}
Choose $\beta\sigma_*>K+2$ and $\beta_*\geq m\beta$, and put
\begin{equation}
 L_S=1+\log(2S),\qquad D=L_S^{\beta_*}e^{B_*}.
 \label{upper:global-cost}
\end{equation}
Then
\begin{align}
 \mathbb E[S^pD^s]&\leq C_p\prod_jF_j,\notag\\
 \mathbb E[(1+N_{\rm pre}+V_{\rm pre}+|\log D|)D^s]
                    &\leq C(1+V_{\rm pre})\prod_jF_j.
 \label{upper:pooled-moments}
\end{align}
If $w_0=X+\log D<0$, choose
\begin{equation}
 h=\left\lceil\log(2S)-w_0\right\rceil.
 \label{upper:global-h}
\end{equation}
At every reached hard owner, whose previous proper outputs continued,
the right-hit gate holds with a fixed dimension constant.
\end{lemma}

\begin{proof}
Expand the separate sums and floors using subadditivity at each
$\sigma_j<1$.  Apply \eqref{upper:subset-moment} to every resulting
owner subset, and sum $(1+h)^{K-\beta\sigma_j}$.  Since
$sr_j=\sigma_j$, this gives the first bound in
\eqref{upper:pooled-moments}.  Use
\eqref{upper:subset-count-moment} for the raw count.  The floors give
$(-\log D)_+\leq C\log(2+V_{\rm pre})$.  The pointwise polynomial
shot bound gives $\log_+D\leq C(1+\log S)$, which is absorbed by a
fixed additional shot moment.  This proves the second bound.

Set $\zeta_\ell=\eta_\ell-X/r_\ell-\theta_\ell$ and
\begin{equation}
 c_j=\theta_j-\sum_{\ell<j}d_\ell\theta_\ell\quad(j<m),
 \qquad c_m=-\sum_{\ell<m}d_\ell\theta_\ell.
 \label{upper:offsets}
\end{equation}
Direct telescoping gives, including $j=m$,
\begin{equation}
 c_j+\sum_{\ell\leq j}\log a_\ell=B_*.
 \label{upper:offset-cancellation}
\end{equation}
The local level at a hard owner satisfies
\begin{align}
 h_{\rm eff}+x_j
 &\leq h+P_{j-1}+x_j\notag\\
 &=h+X+c_j-\sum_{\ell<j}d_\ell\zeta_\ell
 \leq h+X+c_j.
 \label{upper:gate-private-cancellation}
\end{align}
The first inequality holds for both signs of $P_{j-1}$, since
$\min(P_{j-1},0)\leq P_{j-1}$.  The last uses the true previous
continuations $\zeta_\ell>0$.  The choice
\eqref{upper:global-h} yields $h+X+B_*\leq\log(2S)+C$.
Therefore
\begin{equation}
 h_{\rm eff}+x_j+\log\alpha_h-\log S
 \leq C-\sum_{\ell\leq j}\log a_\ell
 \leq C+\sum_{\ell\leq j}D_\ell\log(2+T_\ell).
 \label{upper:gate-past-only}
\end{equation}
All future-owner scales have canceled.  At a hard owner,
\eqref{upper:dominance} implies that every preceding owner's native
length is at most $V_{\rm old}\leq CM$ and the current $T_j\leq CM$.
Thus the last expression is at most $D_{\rm gate}\log(2+M)+C$.
Choose $A_1$ large relative to $D_{\rm gate}$ and $A_0$; direct
comparison of the two logarithmic walls then prevents a right-governed
first hit before $M/2$.  This is the required gate.  No sign assumption
on an incoming output and no bound on a future length were used.
\end{proof}

\begin{remark}[Order of constants]
\label{upper:constant-order}
First fix dimension compactness, rarity and deterministic bank scales.
Choose the prefix floors $D_A$, owner floors $D_j$, and then
$D_{\rm gate}$ from those floors.  Choose $A_1$ next, including the
condition $A_1\sigma_*>2$ for position sums.  Determine the scalar
reserve exponent $K$ after this choice, and finally choose $\beta$
and $\beta_*$.  If a scalar effective-bank threshold depends on $A_1$,
the bounded effective-length case is treated by the probability bound
one and the same middle atom: its persistence ratio is bounded below
by a fixed positive constant.  One does not reduce the original rarity
cutoff in a way which feeds back into $D_{\rm gate}$.  Increasing
$\beta_*$ only improves the gate.
\end{remark}

\begin{theorem}[Actual causal rising profile]
\label{upper:actual-profile}
For the original complete divisor union there is a prefix-uncolored
$D>0$, unread by the final ordinary suffix, such that at its TRUE root
\begin{equation}
 I_{\rm actual}\leq C J_{\beta_*}(X+\log D),\qquad
 J_b(w)=\min\{1,C_b e^w(1+(-w)_+)^b\},
 \label{upper:actual-profile-bound}
\end{equation}
and $D$ satisfies \eqref{upper:pooled-moments}.  All original failure
words are included in this upper bound.
\end{theorem}

\begin{proof}
Use all genuine left-failure covers, keeping their union cost
$CS e^{-h}$.  On the survivors route owners in order.  At a hard owner
allow its internal right stop; otherwise complete its route and stop
at the first $\eta_j\leq\xi_j$.  If there is no proper stop, use the
complete terminal word.  These are actual projected branches.  A
later private filter is never charged at an earlier projection.

For a complete stop at owner $j<m$, the exact capacity exponent is
\begin{equation}
 P_j=X+c_j+\zeta_j-\sum_{\ell<j}d_\ell\zeta_\ell.
 \label{upper:stop-exponent}
\end{equation}
Condition on its true $\mathcal Q_{j-1}$ fiber and apply the negative
kernel, then integrate earlier routes backwards with their positive
kernels.  Each bound is uniform on every genuine preceding fiber.
The result is at most
$Ce^{X+c_j}\prod_{\ell\leq j}b_\ell(h)$.

For an internal hard right stop, its reference exponent is exactly
\[
 P_{j-1}+x_j=X+c_j-\sum_{\ell<j}d_\ell\zeta_\ell.
\]
Use the separate right-stop assertion of
Lemma~\ref{upper:hard-kernels}, followed by the same backward positive
kernels.  It has the same upper bound.  For the terminal branch,
\eqref{upper:carry-identities} gives
$P_m=X+c_m-\sum_{\ell<m}d_\ell\zeta_\ell$.  Integrate its true
native mask first and then use the earlier positive kernels.

The conditional tower integrates each future child word by a uniform
bound before varying the current native count.  The current middle
therefore retains its conditional binomial law, and ancestor atoms
retain their original probabilities.

By \eqref{upper:pool-bounds} and
\eqref{upper:offset-cancellation}, each branch is at most
$Ce^{X+B_*}(1+h)^{j\beta}$.  There are dimension-many branch types,
so
\begin{equation}
 I_{\rm actual}\leq C\min\{1,
          Se^{-h}+e^{X+B_*}(1+h)^{m\beta}\}.
 \label{upper:pre-profile}
\end{equation}
For $w_0<0$, Lemma~\ref{upper:pool-gate} permits the reserve in
\eqref{upper:global-h}.  It satisfies
$1+h\leq C L_S(1+(-w_0)_+)$, and
\eqref{upper:global-cost} now gives
\eqref{upper:actual-profile-bound}.  For $w_0\geq0$, use the global
original rectangle and no gated cover.  The case $m=1$ has no proper
frontier and follows directly from its terminal native masks.
\end{proof}

\subsection{One unread ordinary root and the arithmetic conclusion}
\label{upper:root-section}

We state explicitly the earlier one-root input, including its range.
After the proved whole-original-block trimming needed by terminal
nonzero banks, the unread ordinary suffix has length $U$, largest
block $M_*$, side lengths $A_*,B_*$, and deterministic mean surplus
$\Delta_*$.  It retains the same intrinsic ordinary factor
\begin{equation}
 \mathcal B_I\asymp
 (1+U)^{-1/2}\min\left\{1,
 \frac{(1+B_*)(1+\Delta_*+A_*)}{1+M_*}\right\},
 \qquad \Delta_*\geq cV_{\rm pre}.
 \label{upper:ordinary-factor}
\end{equation}
The trimming leaves the largest root block unchanged and reads no
clock needed by the new proper-owner banks.  On the same original
tuple its genuine endpoint and largest-block ordinary rectangles may
be combined with \eqref{upper:actual-profile-bound}.

\begin{proposition}[The joint ordinary root estimate]
\label{upper:root-contract}
Let $D>0$ and $\kappa$ be arbitrary original-prefix-measurable numbers,
and write $X=\kappa+(\sigma_m/s)Z_{\rm suffix}$.  If the same tuple
has the profile \eqref{upper:actual-profile-bound}, its combination
with the ordinary suffix rectangles satisfies
\begin{equation}
 \mathbb E_{\rm suffix}[e^{-sX}I_{\rm actual}\mid G_{\rm pre}]
 \leq C D^s e^{-\sigma_m H_-}(1+U)^{-1/2}
 \min\left\{1,
 \frac{(1+B_*)(1+H_++A_*)}{1+M_*}\right\},
 \label{upper:root-bound}
\end{equation}
Here $H=(-\log D-\kappa)/(\sigma_m/s)$.
This is valid for every translated root, all original counts, and
the fixed polynomial profile exponent.  It contains the one largest-
block atom jointly with its ballot denominator, with all $n+1$ gaps.
The rectangle's ordered-sum tail is Lemma~\ref{ordsum:rectangle-tail},
including its separate count-zero case.
No additional Poisson atom is to be appended.
\end{proposition}
\begin{proof}
The same-count layer calculation, including every translate and the
polynomial outgoing profile, is given in Section~\ref{updetail:ordinary-root}.
\end{proof}

\begin{theorem}[Upper bound with the same intrinsic factor]
\label{upper:sharp-upper}
At the fixed lifted physical optimizer,
\begin{equation}
 K_k(\ell^+)\leq C_H e^{-J}
               \mathcal C_{\rm count}\mathcal B_I
               \mathcal P_{\rm native}.
 \label{upper:final-upper}
\end{equation}
Together with the previously established complete intrinsic arithmetic
lower, this is a two-sided estimate with that same finite factor.
Actual bounded-parameter stability transfers it to the original vector;
the original multiplication-table and $H$ comparisons are then applied
only on their original admissible domains.
\end{theorem}

\begin{proof}
The profile coefficient of Theorem~\ref{upper:actual-profile} reads no
uncolored point of the unread suffix.  Apply
\eqref{upper:root-bound} ONCE at the actual original root.  The exact
prefix normalization gives
$|\kappa|\leq C(1+V_{\rm pre}+N_{\rm pre})$.
For its first alternative use $e^{-\sigma_mH_-}\leq1$ and the first
bound of \eqref{upper:pooled-moments}.  For its second alternative use
\[
 1+H_++A_*\leq C(1+A_*+V_{\rm pre}+N_{\rm pre}+|\log D|)
\]
and the second bound.  Average these two upper bounds separately and
then take their minimum.  The inequality in
\eqref{upper:ordinary-factor} compares the result with
$C\mathcal B_I\prod_jF_j$.

All actual geometric covers preceded the positive prime-series and
cell enlargement.  The original exponential change of measure bounds
$K_k$ above by $C_H e^{-J}\mathbb E[e^{-sX}I_{\rm actual}]$;
consistently rounded harmonic budgets have only a fixed normalization
error.  Thus
\eqref{upper:factor-product} proves \eqref{upper:final-upper} with
$\mathcal P_{\rm ret}$.  Changing between the finitely many fixed
rarity thresholds only changes a dimension constant, so the full
$\mathcal P_{\rm native}$ is equivalent.  The current correlations
were retained by Proposition~\ref{upper:established-inputs} and the
bank comparisons.  Finally invoke the actual lower and the arithmetic
bounded-lift stability.
\end{proof}
